\chapter{분포 --- 데이터가 따르는 모양}
\label{ch:distributions}
% source: deck 6 slide 1 (내용 없음, 제목만)

\begin{keyidea}
분포를 가정하는 것은 \analogy{투영법}을 고르는 일이다. 편리해지지만 반드시
왜곡이 따르고, 어디가 왜곡되었는지는 검증해야만 알 수 있다.
\end{keyidea}

\begin{inbrief}
  \item[Why] 데이터의 모양을 이름 있는 \term{분포}(Distribution) 하나로
             요약하면, 모수(Parameter) 몇 개만으로 확률을 계산할 수 있다.
  \item[When] 히스토그램의 모양(대칭, 한쪽 꼬리, 평평함, 두 값)이 보이고,
              그 모양을 계산에 쓰고 싶을 때.
  \item[Where] 정규(Normal), 지수(Exponential), 멱법칙(Power Law, 연속),
               균등(Uniform, 이산과 연속), 베르누이(Bernoulli),
               이항(Binomial, 이산). 뒤의 검정(\Cref{ch:hypothesis})과
               정규화(\Cref{ch:qq-normalization})가 이 장의 가정 위에 선다.
  \item[How] 히스토그램 모양으로 후보를 고른다 $\to$ 모수를 추정한다
             $\to$ CDF를 겹쳐 얼마나 어긋나는지 눈으로, 또는 검정으로
             확인한다.
\end{inbrief}

이 장의 출처는 \texttt{notes/6\_DataDistribution.pptx}(30 슬라이드) 하나다.
슬라이드가 외부 문헌을 밝히지 않는 자리는 슬라이드 자체를 출처로 적는다.

% =====================================================================
\section{배경}
\label{sec:dist-background}
% source: deck 6 slide 2, 3

\paragraph{왜 분포인가.} 측정값을 하나씩 늘어놓으면 숫자 더미일 뿐이지만,
같은 값 근처가 얼마나 자주 나오는지를 세어 모양으로 그리면(히스토그램,
\Cref{sec:histogram}) 그 모양에서 규칙이 보인다. 모양이 가장 높이 솟은
자리가 가장 흔한 값이고, 낮은 자리는 드문 값이다. 이 장은 실제로 자주
나타나는 몇 가지 모양에 이름을 붙이고, 그 이름 하나와 모수 몇 개만으로
확률을 계산하는 법을 다룬다.

이 장은 다음 개념을 그대로 가져와 쓴다.

\begin{itemize}
  \item 평균과 표준편차(\Cref{sec:center}, \Cref{def:spread}): 분포의
        중심과 폭을 요약하는 숫자. 이 장의 정규분포, 균등분포, 지수분포
        모두 평균과 분산(또는 표준편차)으로 모양이 정해진다.
  \item 히스토그램과 구간 나누기(\Cref{sec:histogram}): 분포의 모양을
        눈으로 보는 도구.
  \item 왜도와 그 부호(\Cref{sec:skewness}, \Cref{sec:skew-sign}):
        정규분포처럼 대칭이 아닌 모양을 설명하는 말.
\end{itemize}

\term{확률밀도함수}(Probability Density Function, PDF)와
\term{확률질량함수}(Probability Mass Function, PMF)의 구분도 여기서
한 번 깐다. 연속형 변수(어떤 실수든 값이 될 수 있는 변수, 예: 키)는
PDF를 쓰고, 그 값 자체는 확률이 아니라 \textbf{밀도}다. 확률은 PDF
곡선 아래 \textbf{넓이}로만 정해진다. 이산형 변수(정해진 값만 가지는
변수, 예: 주사위 눈)는 PMF를 쓰고, PMF 값 자체가 그 값이 나올 확률이다.

슬라이드 3은 흔한 분포를 연속형과 이산형으로 나눈다. \Cref{tab:dist-map}이
그 목록을 이 장의 절과 짝짓는다. 슬라이드는 정규 $\to$ 균등 $\to$
지수·멱법칙 $\to$ 베르누이 $\to$ 이항 순서로 나열하지만, 이 장은 지수분포와
멱법칙을 먼저 묶고(둘 다 ``꼬리가 얼마나 빨리 줄어드는가''를 다룬다)
균등분포를 뒤로 옮겼다. 슬라이드 순서를 그대로 따르지 않는다.

\begin{table}[ht]\centering
\begin{tabularx}{\linewidth}{@{}l l X@{}}
\toprule
분포 & 연속/이산 & 이 장의 절 \\
\midrule
정규(Normal)     & 연속 & \Cref{sec:normal} \\
지수(Exponential) & 연속 & \Cref{sec:exponential} \\
멱법칙(Power Law) & 연속 & \Cref{sec:powerlaw} \\
균등(Uniform)     & 연속·이산 둘 다 & \Cref{sec:uniform} \\
베르누이(Bernoulli) & 이산 & \Cref{sec:binomial} \\
이항(Binomial)     & 이산 & \Cref{sec:binomial} \\
\bottomrule
\end{tabularx}
\caption{슬라이드 3의 분포 목록과 이 장의 절.}
\label{tab:dist-map}
\end{table}

% =====================================================================
\section{정규분포}
\label{sec:normal}

\begin{inbrief}
  \item[Why] 자연에서 반복해서 나타나는 종 모양의 분포를 평균과 분산,
             모수 둘만으로 완전히 표현하기 위해.
  \item[When] 측정값이 평균 주변에 좌우 대칭으로 몰리고, 평균에서 멀어질수록
              드물어질 때.
  \item[Where] 키, 몸무게, 성적처럼 자연·사회 현상 전반. 이 책에서는
               패스트푸드 \texttt{calories}.
  \item[How] 평균 $\mu$로 중심을, 표준편차 $\sigma$로 폭을 잡는다. PDF
             곡선 아래 넓이가 항상 1이 되도록 높이가 자동으로 맞춰진다.
\end{inbrief}

% source: deck 6 slide 4
정규분포(Normal Distribution)는 \term{가우스분포}(Gaussian Distribution)
또는 \term{종 모양 곡선}(Bell Curve)이라고도 부른다. 자연에서 특히 흔한
이유는 회사의 급여 분포, 키, 몸무게, 성적처럼 많은 독립적인 원인이 더해져
만들어지는 값들이 대체로 이 모양을 따르기 때문이다.

% source: deck 6 slide 5
세로축을 읽을 때 흔히 하는 오해가 하나 있다. PDF의 세로축 값은
확률이 \textbf{아니라} 상대적 가능성(밀도)이다. 아주 크거나 아주 작은
값일수록 곡선이 낮아 가능성이 작고, 평균 근처일수록 곡선이 높아
가능성이 크다는 뜻일 뿐, 특정 $x$ 하나가 나올 확률은 여전히 그 점에서의
넓이(사실상 0)로 정해진다. 실제 확률은 구간에 대해 곡선 아래 넓이를
적분해야 나온다.

% source: deck 6 slide 6
\Cref{sec:dist-background}의 배경 개념(연속 변수는 PDF, 확률은 넓이로만
읽는다)이 이 절 전체의 전제다. 예를 들어 아기 키 분포와 성인 키 분포를
겹쳐 그리면, 두 곡선 모두 각자의 평균을 중심으로 대칭이지만 폭(표준편차)이
다르다. 아기 키의 곡선이 성인 키의 곡선보다 뾰족하고 높은데, 이는 아기가
더 ``특별해서''가 아니라 \textbf{곡선 아래 넓이가 항상 1로 고정}되어
있기 때문이다. 아기 키는 성인 키보다 흩어짐(표준편차)이 작아 폭이
좁으므로, 넓이 1을 유지하려면 높이가 커져야 한다.

% source: deck 6 slide 10 (주석만 추가, 정의 자체는 이전 슬라이드와 동일)
\begin{definition}[정규분포]
\label{def:normal}
\term{정규분포}(Normal Distribution) $X \sim \mathcal{N}(\mu, \sigma^2)$ 의
확률밀도함수(Probability Density Function)는
\[
  f(x) = \frac{1}{\sigma\sqrt{2\pi}}
         \exp\!\left(-\frac{(x-\mu)^2}{2\sigma^2}\right)
\]
\end{definition}

% source: deck 6 slide 7
정규분포의 폭을 $\sigma$ 단위로 재면 어디까지가 ``흔한 범위''인지 규칙이
생긴다. 이를 68-95-99.7 규칙이라 부른다.

\begin{pitfall}[95\%가 아니라 95.45\%]
슬라이드 값은 $\pm 2$ 표준편차 안에 $95\%$, 다시 계산하면 $95.45\%$다. 표준정규분포의
누적분포함수로 $P(|Z|<2) = \Phi(2)-\Phi(-2)$를 계산하면 $0.9545$가
나온다. 정확히 $95\%$가 되는 경계는 $\pm 2$가 아니라 $\pm \mathbf{1.96}$
표준편차다(\Cref{sec:ci}에서 신뢰구간의 $1.96$과 같은 숫자). 슬라이드 7의
그림은 오히려 $1.96\sigma$ 쪽 눈금을 맞게 그려 두었으므로,
틀린 것은 본문의 ``$95\%$''라는 말뿐이다.
\end{pitfall}

\Cref{fig:dist-sigma-bands}가 패스트푸드 \texttt{calories}에 맞춘 정규
곡선 위에 $\pm 1\sigma$, $\pm 2\sigma$ 구간을 표시한다.

\begin{figure}[ht]\centering
  \includegraphics[width=0.85\linewidth]{06_02_normal_sigma_bands.png}
  \caption{패스트푸드 \texttt{calories}에 맞춘 정규 곡선과
  $\pm 1\sigma$($68.27\%$), $\pm 2\sigma$($95.45\%$) 구간. 실제 데이터
  중 $\pm 2\sigma$ 안에 든 비율은 $96.0\%$로, 이론값 $95.45\%$와
  가깝다.}
  \label{fig:dist-sigma-bands}
\end{figure}

\Cref{fig:normal-fit}(\Cref{ch:qq-normalization})은 같은 \texttt{calories}
데이터에 최대우도로 정규 곡선을 맞춘 그림이다. 히스토그램은 종 모양에
가까워 보이지만, 바로 아래 함정과 \Cref{sec:qq-apply}의 Q-Q 플롯은
정규성을 기각한다.

% source: deck 6 slide 8
분포를 이동하거나 척도를 바꾸는 조작도 자주 쓴다. $Z = (X-\mu)/\sigma$로
바꾸면 평균 0, 표준편차 1인 표준정규분포가 되는데, 이 변환의 자세한
성질과 다른 변환(범위 정규화 등)과의 비교는
\Cref{sec:normalization}(\Cref{ch:qq-normalization})에서 다룬다.

% source: deck 6 slide 9
정규분포가 아닌 모양도 있다. 오른쪽으로 긴 꼬리(양의 왜도), 왼쪽으로 긴
꼬리(음의 왜도), 대칭(왜도 0에 가까움), 이 세 모양의 정의와 판정 규칙은
\Cref{sec:skewness}와 \Cref{sec:skew-sign}이 다룬다(\Cref{sec:dist-background}에서
이미 배경으로 소개했다). 이 장은 링크만 걸고 새로 설명하지 않는다.

\begin{pitfall}[정규성 검정의 함정]
같은 표본에서 $\mu, \sigma$ 를 추정한 뒤 그 표본으로 KS 검정(Kolmogorov--Smirnov
Test)을 하면 p-value가 낙관적으로 나온다. 이 경우 Shapiro--Wilk 검정을 쓴다.
% TODO: calories 의 Shapiro p = 0.00038 -- 정규성이 기각된 사례로 정리.
\end{pitfall}

% =====================================================================
\section{지수분포}
\label{sec:exponential}

\begin{inbrief}
  \item[Why] 다음 사건까지 걸리는 시간(대기 시간, 수명)을 모수 하나
             $\lambda$로 표현하기 위해.
  \item[When] 사건이 평균적으로 일정한 비율(rate)로 무작위하게 일어날 때.
  \item[Where] 노트북 수명, 커피숍 손님 도착 간격, 패스트푸드
               \texttt{sugars}.
  \item[How] $\lambda$는 단위 시간당 사건 발생 비율. CDF
             $P(X<x) = 1 - e^{-\lambda x}$로 ``$x$ 안에 사건이 일어날
             확률''을 구한다.
\end{inbrief}

% source: deck 6 slide 15
지수분포(Exponential Distribution)는 사건과 사건 사이에 걸리는 시간을
모형화할 때 쓴다. 예를 들어 가게에서 다음 손님이 올 때까지 기다리는
시간이 그렇다.

\begin{pitfall}[모수 범위]
슬라이드 값은 모수 범위가 ``$0 < \lambda \le 1$''이고, 다시 계산하면
\textbf{$\lambda > 1$도 얼마든지 가능하다}. 올바른 범위는 $\lambda > 0$이다. 슬라이드 17
자신이 시간당 손님 10명, 즉 $\lambda = 10$을 예제로 쓰고 있어 슬라이드
15의 ``$\le 1$'' 제한과 슬라이드 17이 서로 어긋난다.
\end{pitfall}

% source: deck 6 slide 16
\begin{example}[노트북 수명]
\label{ex:dist-laptop}
델 노트북의 평균 수명이 5년일 때, 3년 안에 고장 날 확률을 구한다(슬라이드
16).
\begin{align*}
  \lambda &= 1/5 = 0.2      && \text{평균 5년의 역수}\\
  P(X<3) &= 1 - e^{-\lambda \cdot 3}  && \text{지수분포 CDF}\\
         &= 1 - e^{-0.2\times 3} && \text{$\lambda$ 대입}\\
         &= 1 - e^{-0.6}   && \text{정리}\\
         &= 1 - 0.548812   && \text{$e^{-0.6}$ 계산}\\
         &= \mathbf{0.4512} && \text{슬라이드 값과 확인}
\end{align*}
\end{example}

% source: deck 6 slide 17
\begin{example}[커피숍 손님 도착]
\label{ex:dist-coffee}
시간당 평균 10명이 오는 커피숍(즉 $\lambda = 10$)에서, 다음 손님까지
평균 대기 시간과 5분 안에 손님이 올 확률을 구한다.
\begin{align*}
  \text{평균 대기 시간} &= 1/\lambda           && \text{평균 대기 시간 공식}\\
    &= 1/10\,\text{시간}                        && \text{$\lambda=10$ 대입}\\
    &= 0.1\,\text{시간}                          && \text{소수로 정리}\\
    &= \mathbf{6}\,\text{분}                     && \text{시간을 분으로 환산}
\end{align*}
\end{example}

\begin{pitfall}[슬라이드 값은 0.0801, 다시 계산하면 0.5654]
$5$분을 시간 단위로 바꾸면 $5/60$시간이다.
\begin{align*}
  P(X<5/60) &= 1 - e^{-\lambda \cdot 5/60}   && \text{지수분포 CDF}\\
            &= 1 - e^{-10\times 5/60}        && \text{$\lambda=10$ 대입}\\
            &= 1 - e^{-0.8333}               && \text{지수 정리}\\
            &= 1 - 0.434598                  && \text{$e^{-0.8333}$ 계산}\\
            &= \mathbf{0.5654}               && \text{슬라이드 0.0801과 다르다}
\end{align*}
슬라이드의 $\mathbf{0.0801}$은 $\lambda$를 빠뜨리고
$1-e^{-5/60} = \mathbf{0.0800}$을 계산한 값과 거의 같다(넷째 자리에서만
차이). 상식으로도 검산된다. 평균 대기 시간이 $6$분인 가게에서 $5$분
안에 손님이 올 확률이 $8\%$뿐일 수는 없다. 절반을 넘는 것이 맞다.
\end{pitfall}

\subsection{기하분포로 근사하기}
\label{subsec:geometric}
% source: deck 6 slide 18

\paragraph{고치는 문제.} 지수분포는 연속 시간의 모형이라 ``정확히 몇 분째에
사건이 일어나는가''를 묻기는 어렵다. 슬라이드 18은 이를 1분 단위로 쪼갠
이산 질문으로 바꾼다: 첫 방문자가 1분째, 2분째, 3분째 각각에 올 확률은
얼마인가.

\paragraph{아이디어.} 1분마다 ``이번 1분 동안 방문자가 온다''는 베르누이
시행이 반복된다고 \textbf{근사}한다. 성공 확률은 슬라이드 17과 같은 방식으로
평균 대기 시간의 역수에서 얻는다(1분 평균 도착률 $\lambda = 1/20 = 0.05$,
\Cref{ex:dist-website} 참고). 이때 ``$k$번째 분에 처음 성공한다''는 확률은
기하분포(Geometric Distribution)
\[
  P(X=k) = (1-p)^{k-1} p
\]
를 따른다. 이는 근사이지 \textbf{오류가 아니다}. 실제 도착은 분 중간
어디서나 일어날 수 있는데, 그 안의 정확한 시점을 무시하고 ``그 분 안에서
왔는가/아닌가''만 보기 때문에 지수분포의 정확한 값과 조금 어긋난다.

\begin{example}[웹사이트 첫 방문자]
\label{ex:dist-website}
평균 20분에 한 번 방문자가 온다면 $\lambda = 1/20 = 0.05$(분당)다.
\begin{align*}
  P(0<x<1) &= p = \mathbf{0.05}          && \text{1분째 성공}\\
  P(1<x<2) &= (1-p)\,p                    && \text{1분 놓치고 2분째 성공}\\
           &= 0.95\times 0.05             && \text{$p=0.05$ 대입}\\
           &= \mathbf{0.0475}              && \text{정리}\\
  P(2<x<3) &= (1-p)^2 p                   && \text{2분 놓치고 3분째 성공}\\
           &= 0.95^2\times 0.05           && \text{$p=0.05$ 대입}\\
           &= \mathbf{0.045125}            && \text{정리}
\end{align*}
\end{example}

\Cref{tab:dist-geometric}가 이 기하분포 근사값을 지수분포의 정확한 값과
나란히 놓는다. 지수분포의 정확한 값은
$P(k-1<X<k) = e^{-\lambda(k-1)} - e^{-\lambda k}$로 구한다.

\begin{table}[ht]\centering
\begin{tabularx}{\linewidth}{@{}X c c c@{}}
\toprule
분(1분 구간) & 기하분포 근사 & 지수분포 정확값 & 차이 \\
\midrule
1분째 & 0.05000 & 0.04877 & 0.0012 \\
2분째 & 0.04750 & 0.04639 & 0.0011 \\
3분째 & 0.045125 & 0.04413 & 0.0010 \\
\bottomrule
\end{tabularx}
\caption{첫 방문자 도착 확률: 기하분포 근사값과 지수분포 정확값 비교.
근사와 정확값의 차이는 $0.001$ 안팎으로 작다.}
\label{tab:dist-geometric}
\end{table}

이 근사가 성립하는 이유는 각 1분 구간의 성공 여부가 \textbf{독립}이기
때문이다(\Cref{sec:ht-background}의 곱셈 규칙과 같은 전제). 독립이
아니라면(예: 손님이 몰려오는 시간대) 이 근사는 무너진다.

\Cref{fig:dist-exponential}가 \texttt{sugars} 데이터에 맞춘 지수분포의
PDF와 CDF를 함께 보여준다.

\begin{figure}[ht]\centering
  \includegraphics[width=0.95\linewidth]{06_04_exponential_fit.png}
  \caption{패스트푸드 \texttt{sugars}에 맞춘 지수분포. 평균
  $13.29\,\text{g}$에서 $\lambda = 1/\text{평균} = 0.0753$을 얻는다.}
  \label{fig:dist-exponential}
\end{figure}

% =====================================================================
\section{멱법칙}
\label{sec:powerlaw}
% source: deck 6 slide 19, 20, 21, 22

\begin{inbrief}
  \item[Why] 지수분포보다 훨씬 천천히 줄어드는 꼬리(극단적으로 큰 값이
             지수분포의 예측보다 훨씬 자주 나오는 경우)를 표현하기
             위해.
  \item[When] 소수가 대부분을 차지하는 현상, 즉 상위 몇 개가 전체 양의
              대부분을 차지할 때.
  \item[Where] 도시 인구, 단어 등장 빈도, 부의 분포, HW2의 공항 노선 수.
  \item[How] $f(x) \propto x^{-\alpha}$, $x \ge x_{\min}$. $\alpha$가
             작을수록 꼬리가 두껍다(큰 값이 더 자주 나온다).
\end{inbrief}

\term{멱법칙}(Power Law) 분포는 \term{파레토분포}(Pareto Distribution)라고도
부르며, 소수의 항목이 전체 발생 횟수나 값의 대부분을 차지하는 현상을
설명한다. 이른바 ``80/20 법칙''(파레토 원칙, Pareto Principle)이 대표적인
예다. 예를 들어 사회 전체 부의 약 $80\%$를 인구의 약 $20\%$가 가지고
있다는 관찰이 여기서 나온다.

\begin{definition}[멱법칙 분포]
\label{def:powerlaw}
멱법칙 분포는 $x \ge x_{\min}$에서
\[
  f(x) \propto x^{-\alpha}, \qquad \alpha > 1
\]
형태의 확률밀도함수를 갖는다. $\alpha$는 지수(Power Law Exponent),
$x_{\min}$은 멱법칙이 시작되는 하한이다.
\end{definition}

지수분포와 마찬가지로 값이 클수록 드물지만, 지수분포처럼 지수적으로
줄어들지 않고 \textbf{거듭제곱으로만} 줄어들기 때문에 훨씬 천천히
감소한다. 네트워크의 연결 수, 도시 인구, 단어 사용 빈도처럼 소수가
대부분을 차지하는 데이터에서 흔히 등장한다(슬라이드 20).

\term{여집합 누적분포함수}(Complementary CDF, $P(X \ge x)$)로 그리면
멱법칙을 더 쉽게 확인할 수 있다. 예를 들어 ``$x$번 이상 등장한 단어의
개수''를 세어 $x$에 대해 그리면 멱법칙은 로그-로그(log-log) 축에서
직선에 가깝다(슬라이드 21, University of Michigan의 Mark Newman
강의 노트를 슬라이드 자체가 인용한다).

\Cref{fig:dist-powerlaw}가 멱법칙과 지수분포를 로그-로그 축에서 비교한다
(슬라이드 22, Santa Fe Institute의 Aaron Clauset 강의 노트를 슬라이드가
인용한다). 멱법칙은 직선으로, 지수분포는 아래로 휘는 곡선으로
나타난다.

\begin{figure}[ht]\centering
  \includegraphics[width=0.95\linewidth]{06_06_powerlaw_vs_exponential.png}
  \caption{멱법칙(직선)과 지수분포(아래로 휘는 곡선)를 로그-로그 축에서
  비교. 오른쪽 패널은 패스트푸드 \texttt{sugars}의 여집합 CDF로,
  전체 $126$개 항목 중 로그 축에 표시할 수 없는 $0$ 이하 값을 뺀
  $109$개만 그렸다. 그 $109$개조차 꼬리가 짧아 멱법칙이라 부르기엔
  근거가 부족하다.}
  \label{fig:dist-powerlaw}
\end{figure}

이 데이터셋 안에는 멱법칙을 뚜렷이 보이는 변수가 없다. 실제 멱법칙
데이터와 모수 추정(최대우도, $\hat\alpha = 1 + n/\sum_i \ln(x_i/x_{\min})$)은
HW2의 공항 노선 수 예제(\Cref{sec:hw2-estimators},
\Cref{fig:hw2-airport-powerlaw})를 참고한다.

% =====================================================================
\section{균등분포}
\label{sec:uniform}

\begin{inbrief}
  \item[Why] 모든 결과가 똑같이 나올 가능성을 가질 때를 표현하기 위해.
  \item[When] 정해진 범위 안에서 특정 값을 더 선호할 이유가 없을 때.
  \item[Where] 주사위(이산), 난수 생성기(연속). 이 데이터셋에는 실제로
               균등분포를 따르는 변수가 없다.
  \item[How] 이산이면 각 결과에 $1/n$을, 연속이면 구간 폭의 역수를
             밀도로 준다.
\end{inbrief}

% source: deck 6 slide 11
균등분포(Uniform Distribution)는 결과가 정해진 범위 안에서 임의로
나오는 실험을 설명하며, 모든 결과가 똑같이 나올 가능성(Equally
Likely)을 갖는다. 공정한 6면체 주사위를 굴리는 것이 대표적인 이산
균등분포(Discrete Uniform Distribution)의 예로, $P(X=1) = \cdots =
P(X=6) = 1/6$이고 모든 확률의 합은 1이다.

% source: deck 6 slide 12
\begin{example}[연속 균등분포]
\label{ex:dist-uniform-continuous}
구간 $[0, 5]$ 위의 연속 균등분포를 생각한다. 밀도 함수는 구간
전체에서 상수 $h$이고, 넓이가 1이어야 하므로
\begin{align*}
  5 \times h &= 1  && \text{구간 폭 5, 넓이 1}\\
  h &= 1/5          && \text{밀도의 높이}
\end{align*}
슬라이드는 이 식을 ``$5\times X = 1 \to X = 1/5$''로 적어 구하는
대상을 확률변수 $X$처럼 보이게 썼지만, 실제로 구하는 것은 $X$가 아니라
밀도의 \textbf{높이} $h$다. 이 높이 $h=1/5$를 써서 구간 확률을 구한다.
\begin{align*}
  P(1<X<2) &= 1 \times \tfrac15 = \tfrac15
    && \text{폭 1 구간}\\
  P(4<X<4\tfrac13) &= \tfrac13 \times \tfrac15 = \mathbf{1/15}
    && \text{폭 $1/3$ 구간}\\
  P(X=4) &= 0
    && \text{점 하나의 넓이는 0}
\end{align*}
\end{example}

% source: deck 6 slide 14
일반적으로 구간 $[a, b]$ 위의 균등분포는(기호는 \Cref{sec:qq-background}의
함정에서 정한 것과 같이 $[a,b]$로 통일한다) 다음 평균과 표준편차를
갖는다.
\[
  \E[X] = \frac{a+b}{2}, \qquad
  \sd(X) = \frac{b-a}{\sqrt{12}}
\]
$[0,5]$ 구간에 대입하면
\begin{align*}
  \E[X] &= (0+5)/2 = 2.5            && \text{평균}\\
  \sd(X) &= 5/\sqrt{12} = 1.4434    && \text{표준편차}
\end{align*}

% source: deck 6 slide 13
슬라이드 13은 생후 8주 아기의 웃는 시간 55개를 난수로 만드는 예제를
든다. 균등분포에서 난수를 생성하는 방식(각 값이 똑같이 나올 가능성을
갖도록)을 보여주는 자리이며, 이 책은 슬라이드를 출처로만 인용하고 개별
값은 옮겨 적지 않는다.

\Cref{fig:dist-uniform}이 이산 균등분포(주사위)와 연속 균등분포
$[0,5]$를 함께 보여준다.

\begin{figure}[ht]\centering
  \includegraphics[width=0.95\linewidth]{06_05_uniform_reference.png}
  \caption{이산 균등분포(주사위, 왼쪽)와 연속 균등분포 $[0,5]$(오른쪽).
  이 데이터셋 안에는 균등분포를 따르는 변수가 없어 이론적 참고 그림으로
  그렸다.}
  \label{fig:dist-uniform}
\end{figure}

\begin{pitfall}[serving\_size는 균등분포가 아니다]
\texttt{serving\_size}가 균등분포인지
경험적 CDF와의 최대 세로 격차(\Cref{sec:cdf}의 $D$)로 확인하면
\textbf{0.152 대 0.101}이다. 균등분포로 맞췄을 때 최대격차가 $0.152$로,
정규분포로 맞췄을 때의 $0.101$보다 크다. 즉 \texttt{serving\_size}는
균등분포보다 정규분포에 더 가깝다.
\end{pitfall}

% =====================================================================
\section{베르누이와 이항분포}
\label{sec:binomial}

\begin{inbrief}
  \item[Why] 결과가 성공/실패 둘뿐인 시행 하나(베르누이), 또는 그런
             시행을 여러 번 반복했을 때의 성공 횟수(이항)를 표현하기
             위해.
  \item[When] 이진 결과(예/아니오, 성공/실패)를 셀 때.
  \item[Where] 동전 던지기, 지지/반대 설문, 패스트푸드
               \texttt{sodium}이 기준치를 넘는지 여부.
  \item[How] 베르누이는 성공확률 $p$ 하나. 이항은 여기에 시행 횟수
             $n$을 더해 $n$번 중 성공 횟수의 분포를 구한다.
\end{inbrief}

% source: deck 6 slide 23
베르누이분포(Bernoulli Distribution)는 각 데이터를 이진값(Binary Value)
하나로 매핑할 수 있을 때 쓴다. 실험 결과가 성공 또는 실패뿐이면, 성공(값
1)에는 확률 $p$를, 실패(값 0)에는 확률 $1-p$를 준다. 동전의 앞/뒤,
찬성/반대가 그 예다.

% source: deck 6 slide 24
\begin{example}[지지율 설문]
\label{ex:dist-favorability}
설문 응답자의 $60\%$가 찬성($1$), $40\%$가 반대($0$)라면 평균과 분산은
다음과 같다.
\begin{align*}
  \mu &= 0.4\times 0 + 0.6\times 1 = \mathbf{0.6}
    && \text{평균, 그러나 분포가 실제로 가질 수 있는 값은 0과 1뿐}\\
  \sigma^2 &= 0.4(0-0.6)^2 + 0.6(1-0.6)^2 && \text{분산: 편차제곱의 기댓값}\\
           &= 0.144 + 0.096 = \mathbf{0.24} && \text{정리}
\end{align*}
\end{example}

% source: deck 6 slide 25
일반적으로 성공확률 $p$인 베르누이분포의 평균과 표준편차는
\[
  \E[X] = p, \qquad \sd(X) = \sqrt{p(1-p)}
\]
이고, 확률질량함수는 $P(X=1)=p$, $P(X=0)=1-p$다(슬라이드 25의
표기와 같다).

% source: deck 6 slide 26
이항분포(Binomial Distribution)는 베르누이 시행을 서로 \textbf{독립}으로
$n$번 반복했을 때, 성공 횟수 $X$의 분포다. 예를 들어 동전을 5번 던져
``뒷면'' 2번이 나올 확률을 구하는 문제가 여기 속한다.

% source: deck 6 slide 27
\begin{example}[동전 5번 중 순서까지 정해진 경우]
\label{ex:dist-coin-order}
동전을 5번 던져 앞의 3번은 실패(성공=뒷면=1, 실패=앞면=0), 뒤의 2번은
성공인 \textbf{정해진 순서 하나}가 나올 확률은
\[
  0.5^5 = \mathbf{0.03125}
\]
다(각 시행이 독립이므로 곱한다, \Cref{sec:ht-background}).
\end{example}

% source: deck 6 slide 28
\begin{example}[순서를 가리지 않는 경우]
\label{ex:dist-coin-combination}
성공이 정확히 2번 나오는 \textbf{모든} 순서를 세려면 조합
$\binom{5}{2}$을 곱한다.
\begin{align*}
  \binom{5}{2} &= 10                      && \text{5번 중 2번을 고르는 방법 수}\\
  P(X=2) &= \binom{5}{2}\times 0.5^5      && \text{순서 무관 확률}\\
         &= 10/32                          && \text{값 대입}\\
         &= \mathbf{0.3125}                && \text{정리}
\end{align*}
\end{example}

% source: deck 6 slide 29
일반적으로 $n$번의 독립 베르누이 시행에서 성공확률 $p$인 이항분포의
평균과 표준편차는
\[
  \E[X] = np, \qquad \sd(X) = \sqrt{np(1-p)}
\]
다(슬라이드 29의 표기와 같다).

\Cref{fig:dist-bernoulli-binomial}이 패스트푸드 \texttt{sodium}이
$1000\,\text{mg}$을 넘는지 여부를 베르누이 시행으로, $n=10$번 뽑았을 때
고나트륨 항목 개수를 이항분포로 보여준다.

\begin{figure}[ht]\centering
  \includegraphics[width=0.95\linewidth]{06_07_bernoulli_binomial.png}
  \caption{패스트푸드 \texttt{sodium} $>1000\,\text{mg}$ 여부의
  베르누이 PMF($p=0.4524$)와, $n=10$번 뽑았을 때 고나트륨 개수의
  이항 PMF(평균 $np=4.52$, 표준편차 $1.57$).}
  \label{fig:dist-bernoulli-binomial}
\end{figure}

% =====================================================================
\section{누적분포함수로 비교하기}
\label{sec:cdf}
% source: deck 6 slide 16-17 (CDF as tool); figure from hw1 data

\begin{inbrief}
  \item[Why] 히스토그램 없이도, 분포 전체를 한 곡선으로 겹쳐 비교하기
             위해.
  \item[When] 데이터가 특정 분포를 따르는지 눈으로 먼저 확인하고 싶을 때.
  \item[Where] 지수분포 CDF로 대기 시간을 구했던 계산(슬라이드 16, 17)이
               이미 CDF를 도구로 쓴 예다.
  \item[How] 경험적 CDF(계단 모양)와 이론 CDF(매끄러운 곡선)를 같은
             축에 겹치고, 가장 크게 벌어지는 지점을 본다.
\end{inbrief}

값 $x$에서의 \term{누적분포함수}(Cumulative Distribution Function, CDF)는
한 줄로 요약하면 $F(x) = P(X \le x)$, 즉 ``$x$ 이하가 전체에서 차지하는
비율''이다. 데이터로 만든 계단 모양의 CDF를
\term{경험적 누적분포함수}(Empirical CDF, $\hat F$)라 부르며, 정식
정의는 \Cref{def:ecdf}(\Cref{ch:compare-dist})가 맡는다. 이 장은 CDF를
지수분포 예제(\Cref{ex:dist-laptop}, \Cref{ex:dist-coffee})에서 이미
계산 도구로 썼으므로, 정의를 다시 쓰지 않고 두 곡선을 겹쳐 보는 용도로만
쓴다.
모집단 CDF의 정식 정의(단조 증가, 우연속, 극한 0과 1)는 \Cref{def:cdf}에 있다.

\Cref{fig:dist-cdf}가 패스트푸드 \texttt{calories}의 경험적 CDF와
정규분포로 맞춘 CDF를 겹쳐 보여준다.

\begin{figure}[ht]\centering
  \includegraphics[width=0.85\linewidth]{06_03_cdf.png}
  \caption{패스트푸드 \texttt{calories}의 경험적 CDF(계단)와 최대우도로
  맞춘 정규 CDF($\hat\mu=532.49$, $\hat\sigma=249.85$). 두 곡선의 최대
  세로 격차 $D=0.1007$이 표시되어 있다.}
  \label{fig:dist-cdf}
\end{figure}

두 곡선이 가장 크게 벌어지는 세로 거리가
\term{콜모고로프-스미르노프 통계량}(Kolmogorov-Smirnov Statistic) $D$다.
정식 정의와 $D$를 이용한 검정 절차는 \Cref{def:ks-statistic}과
\Cref{sec:cd-ks}(\Cref{ch:compare-dist})가 다룬다. 같은 표본에서 모수를
추정한 뒤 이 표본으로 KS 검정을 하면 p값이 낙관적으로 나오므로
(\Cref{sec:normal}의 함정), 이 장에서는 $D$를 \textbf{눈으로 보는} 도구로만
쓰고 p값은 계산하지 않는다. 히스토그램 대신 CDF로 비교하면 구간 폭을
정할 필요가 없다는 장점은 \Cref{sec:dist-check-cdf}(\Cref{ch:qq-normalization})가
더 다룬다.

% =====================================================================
\section{요약}

\Cref{tab:dist-summary}가 이 장에서 다룬 분포를 모수, 평균, 분산과 함께
정리한다.

\begin{table}[ht]\centering
\begin{tabularx}{\linewidth}{@{}l l X c c c@{}}
\toprule
분포(Distribution) & 연속/이산 & 모수 & 평균 & 분산 & 절 \\
\midrule
정규(Normal)     & 연속 & $\mu, \sigma^2$ & $\mu$ & $\sigma^2$
  & \Cref{sec:normal} \\
지수(Exponential)     & 연속 & $\lambda$ & $1/\lambda$ & $1/\lambda^2$
  & \Cref{sec:exponential} \\
멱법칙(Power Law)   & 연속 & $\alpha, x_{\min}$ & (조건부 존재) & (조건부 존재)
  & \Cref{sec:powerlaw} \\
균등(Uniform)     & 연속 & $a, b$ & $(a+b)/2$ & $(b-a)^2/12$
  & \Cref{sec:uniform} \\
베르누이(Bernoulli) & 이산 & $p$ & $p$ & $p(1-p)$
  & \Cref{sec:binomial} \\
이항(Binomial)     & 이산 & $n, p$ & $np$ & $np(1-p)$
  & \Cref{sec:binomial} \\
\bottomrule
\end{tabularx}
\caption{이 장에서 다룬 분포의 모수와 평균, 분산.}
\label{tab:dist-summary}
\end{table}

% =====================================================================
\section{어떻게 적용하는가}
\label{sec:dist-apply}
% source: deck 6 slide 30 (내용 없음, Kahoot 퀴즈)

이 책 안에서 이 장의 분포가 실제로 쓰이는 자리는 다음과 같다.

\begin{itemize}
  \item 패스트푸드 \texttt{calories}: 히스토그램은 종 모양이지만 정규성이
        기각된다(\Cref{fig:normal-fit}, Shapiro-Wilk $p=0.00038$,
        \Cref{sec:normal}의 함정).
  \item 패스트푸드 \texttt{sugars}: 지수분포로 맞추면
        $\hat\lambda = 0.0753$(\Cref{fig:dist-exponential}).
  \item 패스트푸드 \texttt{sodium}: $1000\,\text{mg}$ 초과 여부를
        베르누이/이항으로 요약한다(\Cref{fig:dist-bernoulli-binomial}).
  \item HW2의 공항 노선 수, 영화 평점: 네 후보 분포(멱법칙, 지수, 균등,
        정규)의 모수를 추정하고 적합도를 판정한다
        (\Cref{sec:hw2-estimators}, \Cref{fig:hw2-airport-powerlaw}).
  \item 후보를 고른 뒤에는 Q-Q 플롯으로 한 번 더 확인한다
        (\Cref{sec:qq}, \Cref{fig:qq-calories}).
\end{itemize}

% =====================================================================
\section{앞으로}
\label{sec:dist-future}

\begin{itemize}
  \item \textbf{모수 추정과 Q-Q 플롯.} 이 장은 슬라이드의 공식으로
        모수를 구했다. 최대우도추정(Maximum Likelihood Estimation, MLE)의
        일반적인 절차와 Q-Q로 적합을
        확인하는 법은 \Cref{ch:qq-normalization}이 다룬다.
  \item \textbf{이산분포 위의 정확 검정.} 이항분포에서 관측값이 얼마나
        극단적인지 정확히 계산하는 법은 \Cref{sec:ht-exact}와
        \Cref{sec:pv-discrete}가 다룬다.
  \item \textbf{정규분포 위의 신뢰구간.} $\pm 1.96\sigma$가 왜 $95\%$의
        경계인지, 신뢰구간에 어떻게 쓰이는지는 \Cref{sec:ci}가 다룬다.
  \item \textbf{두 분포를 통계적으로 비교하기.} 이 장은 CDF를 눈으로만
        겹쳐 봤다. 숫자 하나(KS 통계량 $D$)로 ``얼마나 다른가''를
        검정하는 절차는 \Cref{sec:cd-ks}가 다룬다.
\end{itemize}
